\section{Related Work}\label{sec:related}

\subsection{Latent Factor Models and Deep Asset Pricing}

Classical factor models originate from \citet{Ross1976apt}'s Arbitrage Pricing Theory, which establishes that expected returns are linearly related to exposures to a small number of common factors. The empirical implementations by \citet{FamaFrench2015fivefactor} and \citet{Carhart1997persistence} construct tradable factors via characteristic sorts, but are designed for equities and have limited explanatory power for bonds, REITs, and emerging market assets. Statistical approaches such as PCA \citep{BaiNg2002factors} and RP-PCA \citep{LettauPelger2020factors} recover latent factors that maximize explained variation or penalize pricing errors, while IPCA \citep{KellyPruittSu2019ipca} instruments loadings with observable characteristics, enabling time variation within a linear framework. More recently, \citet{GuKellyXiu2021autoencoder} introduced a conditional autoencoder that maps characteristics to factor loadings via neural networks, and subsequent work has extended this paradigm: \citet{ChenPelgerZhu2024deep} adopt adversarial GAN-based SDF estimation with LSTM-extracted macro states, \citet{Epstein2025attention} apply attention mechanisms to construct conditional arbitrage factors, and \citet{Kelly2024large} develop a Portfolio Tangent Kernel decomposition for scaling deep factor models. Unlike these works, which operate on individual stock returns and produce abstract latent factors that cannot be directly traded, our model operates on mutual fund returns and constrains each factor to be the return of an explicit, tradable portfolio via the portfolio bottleneck.

\subsection{Autoencoder and VAE Extensions}

A growing body of work extends the conditional autoencoder architecture along several dimensions. \citet{Duan2022factorvae} introduce FactorVAE, a variational framework that models latent factors as distributions rather than point estimates, enabling uncertainty quantification in noisy financial data. \citet{WangGuo2024rvrae} extend this with recurrent encoders to capture temporal factor dynamics, while \citet{WangSingh2024kan} replace standard MLP layers with Kolmogorov-Arnold networks for improved interpretability of learned loadings. \citet{Gopal2024neuralfactors} propose NeuralFactors, which jointly learns factor exposures and factor returns within a VAE framework, and \citet{Engel2025scaling} scale conditional autoencoders to high-dimensional factor spaces with uncertainty-driven pruning. \citet{Jang2025consensus} take a different approach entirely, using analyst consensus forecasts as a structural bottleneck to inject domain knowledge into the latent space. Unlike these extensions, which focus on probabilistic modeling, alternative network architectures, or information-theoretic bottlenecks---all applied to individual stock returns---our bottleneck is a \emph{portfolio return}: a fundamentally different inductive bias that enforces tradability by construction and operates at the fund level with fund-specific conditioning variables.

\subsection{Portfolio Construction and End-to-End Learning}

The idea that factors should correspond to tradable portfolios dates to \citet{Huberman1987mimicking}, who showed that any factor can be replicated by a portfolio that is maximally correlated with it. \citet{Brandt2009parametric} formalize parametric portfolio policies where weights are linear functions of asset characteristics, enabling direct optimization of portfolio objectives. A parallel thread in decision-focused learning differentiates through optimization layers: \citet{Amos2017optnet} embed quadratic programs as differentiable network layers, \citet{Donti2017taskbased} propose task-based end-to-end learning that minimizes downstream decision loss, and \citet{Zhang2020portfolio} optimize Sharpe ratios directly via differentiable portfolio construction. Unlike these approaches, which learn portfolio allocations to maximize trading performance, our model learns factor-mimicking portfolios that maximize cross-sectional explanatory power---bridging factor modeling and portfolio construction in a single architecture.

\subsection{Machine Learning for Mutual Funds}

\citet{Sharpe1992style} introduced returns-based style analysis, decomposing fund returns into exposures to asset class indices via constrained regression. Recent work applies modern ML to fund analysis: \citet{Kaniel2023skill} train neural networks on stock, fund, and macroeconomic characteristics to identify skilled fund managers, while \citet{DeMiguel2023fund} demonstrate that ML methods exploiting a large set of fund characteristics can generate significant out-of-sample alpha. Unlike these approaches, which use fund characteristics to predict fund performance or select funds, we use fund-level characteristics derived from regulatory filings---sector holdings and asset class allocations---as \emph{conditioning variables for factor construction}. This is a fundamentally different task: rather than predicting which funds will outperform, we learn a latent factor structure that explains \emph{why} fund returns co-move across asset classes.

\paragraph{Positioning.} To our knowledge, no prior work combines (i)~a portfolio bottleneck that enforces factor tradability by construction, (ii)~conditioning on fund-level alternative data from regulatory filings, and (iii)~multi-asset training across equities, bonds, REITs, and emerging markets within a unified autoencoder framework. Our work bridges deep latent factor models (\S\ref{sec:related}.1--2) with portfolio-based factor construction (\S\ref{sec:related}.3), applied to the mutual fund universe (\S\ref{sec:related}.4).
